% GENERATED FILE. Edit canonical lessons, then run npm run generate:books.
% canonical-source-hash: fnv1a64:be01268bcd9e55d9
% canonical-lessons: JA-C348-kesseki-suru, JA-C348-osoi, JA-C348-girigiri, JA-C348-jikoku, JA-C348-renkyuu

\chapter{If You Cannot Come}
\label{ch:ja-if-you-cannot-come}
\hlchaptermodality{348}

\begin{chapteropening}
\textbf{By the end of this chapter:} \emph{I can say what to do if someone is late or absent.}

Complete the last lesson of chapter 348: I can say what to do if someone is late or absent.
\end{chapteropening}

\section[\texorpdfstring{\ja{けっせきする} (kesseki suru)}{けっせきする (kesseki suru)}]{\ja{けっせきする} (kesseki suru) — to be absent; to stay away}

\label{lesson:JA-C348-kesseki-suru}



\begin{quote}
Before the new one: say the Japanese for while you are at it, then the Japanese for a caller.
\end{quote}

\subsection*{You'll want to know: \ja{けっせきする}}
\textbf{\ja{けっせきする}} — \emph{kesseki suru} — \textquotedblleft{}to be absent; to stay away\textquotedblright{}.

\begin{quote}
\textbf{\ja{ねつが} \ja{あったら、がっこうを} \ja{けっせきして} \ja{ください。}} — \emph{netsu ga attara, gakkō o kesseki shite kudasai} — \textquotedblleft{}If you have a fever, stay away from school.\textquotedblright{}
\end{quote}

Polite \textbf{\ja{けっせきします}}, \emph{kessekishimasu}; a \ja{する} verb. Past (polite): \textbf{\ja{むすこが} \ja{がっこうを} \ja{けっせきしました。}}, \emph{musuko ga gakkō o kesseki shimashita}.

\begin{cousinweb}
\textbf{\ja{けっせきする}} (\emph{kesseki suru}) is built from a Sino-Japanese noun and the native verb \emph{suru}, \textquotedblleft{}to do\textquotedblright{}. The noun, \emph{kesseki}, is two Chinese characters read with their on'yomi, readings borrowed from Chinese: \emph{ketsu}, \textquotedblleft{}lack\textquotedblright{} and \emph{seki}, \textquotedblleft{}seat\textquotedblright{} (Middle Chinese \emph{zjek}) (KANJIDIC2; JMdict; Middle Chinese in Baxter's transcription). The two join as \emph{kes-seki}.
\end{cousinweb}

\begin{grammarlens}[title={What you've built}]
One more word for saying what happens if.
\end{grammarlens}

\subsection*{Your turn}
\begin{itemize}
\raggedright
  \item \emph{Say it:} \emph{kesseki suru}
  \item \emph{Say it:} \emph{kesseki suru}, once more
  \item \emph{Say it:} \emph{netsu ga attara, gakkō o kesseki shite kudasai} — If you have a fever, stay away from school.
  \item \emph{Recall:} say the Japanese for while you are at it, then the Japanese for a caller, then say \emph{kesseki suru} again
\end{itemize}

\begin{tcolorbox}[breakable,colback=teal!4,colframe=teal!35!black,title={Before you move on}]
What does \ja{けっせきする} mean? (\textquotedblleft{}To be absent; to stay away\textquotedblright{}.) Say it once more.
\end{tcolorbox}
\section[\texorpdfstring{\ja{おそい} (osoi)}{おそい (osoi)}]{\ja{おそい} (osoi) — slow; late}

\label{lesson:JA-C348-osoi}



\begin{quote}
Before the new one: say the Japanese for a caller, then the Japanese for to be absent.
\end{quote}

\subsection*{You'll want to know: \ja{おそい}}
\textbf{\ja{おそい}} — \emph{osoi} — \textquotedblleft{}slow; late\textquotedblright{}.

\begin{quote}
\textbf{\ja{おそく} \ja{なったら、でんわして} \ja{ください。}} — \emph{osoku nattara, denwa shite kudasai} — \textquotedblleft{}If it gets late, please phone.\textquotedblright{}
\end{quote}

An i-adjective.

\begin{cousinweb}
\textbf{\ja{おそい}} (\emph{osoi}) is a native Japanese word (wago), inherited rather than borrowed from Chinese, and its older origin is not settled. Its kanji was attached to it later as a kun'yomi, a native reading given to a Chinese character (JMdict, KANJIDIC2).
\end{cousinweb}

\begin{grammarlens}[title={What you've built}]
One more word for saying what happens if.
\end{grammarlens}

\subsection*{Your turn}
\begin{itemize}
\raggedright
  \item \emph{Say it:} \emph{osoi}
  \item \emph{Say it:} \emph{osoi}, once more
  \item \emph{Say it:} \emph{osoku nattara, denwa shite kudasai} — If it gets late, please phone.
  \item \emph{Recall:} say the Japanese for a caller, then the Japanese for to be absent, then say \emph{osoi} again
\end{itemize}

\begin{tcolorbox}[breakable,colback=teal!4,colframe=teal!35!black,title={Before you move on}]
What does \ja{おそい} mean? (\textquotedblleft{}Slow; late\textquotedblright{}.) Say it once more.
\end{tcolorbox}
\section[\texorpdfstring{\ja{ぎりぎり} (girigiri)}{ぎりぎり (girigiri)}]{\ja{ぎりぎり} (girigiri) — at the last minute; barely}

\label{lesson:JA-C348-girigiri}



\begin{quote}
Before the new one: say the Japanese for to be absent, then the Japanese for slow.
\end{quote}

\subsection*{You'll want to know: \ja{ぎりぎり}}
\textbf{\ja{ぎりぎり}} — \emph{girigiri} — \textquotedblleft{}at the last minute; barely\textquotedblright{}.

\begin{quote}
\textbf{\ja{いま} \ja{でたら、ぎりぎり} \ja{まにあうでしょう。}} — \emph{ima detara, girigiri maniau deshō} — \textquotedblleft{}If you leave now, you will barely make it.\textquotedblright{}
\end{quote}

\begin{cousinweb}
\textbf{\ja{ぎりぎり}} (\emph{girigiri}) is a native mimetic word, one of the Japanese words that picture a manner or a feeling in sound; JMdict marks it as mimetic. It is inherited, not borrowed, and its older origin is not settled (JMdict).
\end{cousinweb}

\begin{grammarlens}[title={What you've built}]
One more word for saying what happens if.
\end{grammarlens}

\subsection*{Your turn}
\begin{itemize}
\raggedright
  \item \emph{Say it:} \emph{girigiri}
  \item \emph{Say it:} \emph{girigiri}, once more
  \item \emph{Say it:} \emph{ima detara, girigiri maniau deshō} — If you leave now, you will barely make it.
  \item \emph{Recall:} say the Japanese for to be absent, then the Japanese for slow, then say \emph{girigiri} again
\end{itemize}

\begin{tcolorbox}[breakable,colback=teal!4,colframe=teal!35!black,title={Before you move on}]
What does \ja{ぎりぎり} mean? (\textquotedblleft{}At the last minute; barely\textquotedblright{}.) Say it once more.
\end{tcolorbox}
\section[\texorpdfstring{\ja{じこく} (jikoku)}{じこく (jikoku)}]{\ja{じこく} (jikoku) — the hour; the time of day}

\label{lesson:JA-C348-jikoku}



\begin{quote}
Before the new one: say the Japanese for slow, then the Japanese for at the last minute.
\end{quote}

\subsection*{You'll want to know: \ja{じこく}}
\textbf{\ja{じこく}} — \emph{jikoku} — \textquotedblleft{}the hour; the time of day\textquotedblright{}.

\begin{quote}
\textbf{\ja{でんしゃの} \ja{じこくが} \ja{かわったら、おしえて} \ja{ください。}} — \emph{densha no jikoku ga kawattara, oshiete kudasai} — \textquotedblleft{}If the train times change, please tell me.\textquotedblright{}
\end{quote}

\begin{cousinweb}
\textbf{\ja{じこく}} (\emph{jikoku}) is Sino-Japanese (kango): two Chinese characters read with their on'yomi, readings borrowed from Chinese. They are \emph{ji}, \textquotedblleft{}time\textquotedblright{} (Middle Chinese \emph{dzyi}) and \emph{koku}, \textquotedblleft{}time\textquotedblright{} (\emph{khok}) (KANJIDIC2; Middle Chinese in Baxter's transcription).
\end{cousinweb}

\begin{grammarlens}[title={What you've built}]
One more word for saying what happens if.
\end{grammarlens}

\subsection*{Your turn}
\begin{itemize}
\raggedright
  \item \emph{Say it:} \emph{jikoku}
  \item \emph{Say it:} \emph{jikoku}, once more
  \item \emph{Say it:} \emph{densha no jikoku ga kawattara, oshiete kudasai} — If the train times change, please tell me.
  \item \emph{Recall:} say the Japanese for slow, then the Japanese for at the last minute, then say \emph{jikoku} again
\end{itemize}

\begin{tcolorbox}[breakable,colback=teal!4,colframe=teal!35!black,title={Before you move on}]
What does \ja{じこく} mean? (\textquotedblleft{}The hour; the time of day\textquotedblright{}.) Say it once more.
\end{tcolorbox}
\section[\texorpdfstring{\ja{れんきゅう} (renkyū)}{れんきゅう (renkyū)}]{\ja{れんきゅう} (renkyū) — consecutive holidays}

\label{lesson:JA-C348-renkyuu}



\begin{quote}
Before the new one: say the Japanese for at the last minute, then the Japanese for the hour.
\end{quote}

\subsection*{You'll want to know: \ja{れんきゅう}}
\textbf{\ja{れんきゅう}} — \emph{renkyū} — \textquotedblleft{}consecutive holidays\textquotedblright{}.

\begin{quote}
\textbf{\ja{れんきゅうに} \ja{なったら、いなかに} \ja{かえります。}} — \emph{renkyū ni nattara, inaka ni kaerimasu} — \textquotedblleft{}When the long weekend comes, I will go home to the country.\textquotedblright{}
\end{quote}

\begin{cousinweb}
\textbf{\ja{れんきゅう}} (\emph{renkyū}) is Sino-Japanese (kango): two Chinese characters read with their on'yomi, readings borrowed from Chinese. They are \emph{ren}, \textquotedblleft{}connect\textquotedblright{} (Middle Chinese \emph{ljen}) and \emph{kyū}, \textquotedblleft{}rest\textquotedblright{} (\emph{xjuw}) (KANJIDIC2; Middle Chinese in Baxter's transcription).
\end{cousinweb}

\begin{grammarlens}[title={What you've built}]
One more word for saying what happens if.
\end{grammarlens}

\subsection*{Your turn}
\begin{itemize}
\raggedright
  \item \emph{Say it:} \emph{renkyū}
  \item \emph{Say it:} \emph{renkyū}, once more
  \item \emph{Say it:} \emph{renkyū ni nattara, inaka ni kaerimasu} — When the long weekend comes, I will go home to the country.
  \item \emph{Recall:} say the Japanese for at the last minute, then the Japanese for the hour, then say \emph{renkyū} again
\end{itemize}

\begin{tcolorbox}[breakable,colback=teal!4,colframe=teal!35!black,title={Before you move on}]
What does \ja{れんきゅう} mean? (\textquotedblleft{}Consecutive holidays\textquotedblright{}.) Say it once more.
\end{tcolorbox}
